% Part: normal-modal-logic
% Chapter: completeness
% Section: modalities-ccs

\documentclass[../../../include/open-logic-section]{subfiles}

\begin{document}

\olfileid[pt]{nml}{com}{mod}

\olsection{Modalidades e conjuntos consistentes completos}

\begin{explain}
Quando construirmos um modelo~$\mModel{M^\Sigma}$ cujo conjunto de
mundos é dado pelos conjuntos $\Sigma$-consistentes completos~$\Delta$
de alguma lógica modal normal~$\Sigma$, também precisaremos definir uma
relação de acessibilidade~$R^\Sigma$ entre esses ``mundos''. Queremos
que a relação de acessibilidade (e a atribuição $V^\Sigma)$ sejam
definidas de modo que $\mSat{M^\Sigma}{!A}[\Delta]$ se e somente se
$!A \in \Delta$. Como devemos fazer isso?

Uma vez definida a relação de acessibilidade, a definição de verdade
em um mundo garante que \iftag{prvBox}{$\mSat{M^\Sigma}{\Box!A}[\Delta]$
  se e somente se $\mSat{M^\Sigma}{!A}[\Delta']$ para todo $\Delta'$ tal que
  $R^\Sigma\Delta\Delta'$}{$\mSat{M^\Sigma}{\Diamond!A}[\Delta]$ se e somente se
  $\mSat{M^\Sigma}{!A}[\Delta']$ para pelo menos um $\Delta'$ tal que
  $R^\Sigma\Delta\Delta'$}. A demonstração de que
$\mSat{M^\Sigma}{!A}[\Delta]$ se e somente se $!A \in \Delta$ exige que
isso seja verdadeiro, em particular, para !!{formula}s que começam com
um operador modal, O.s.\@,
\iftag{prvBox}{$\mSat{M^\Sigma}{\Box!A}[\Delta]$ se e somente se $\Box !A
  \in \Delta$}{$\mSat{M^\Sigma}{\Diamond!A}[\Delta]$ se e somente se $\Diamond !A
  \in \Delta$}. Combinar essa exigência com a definição de verdade em
um mundo para \iftag{prvBox}{$\Box !A$}{$\Diamond !A$} produz:
\[
\iftag{prvBox}{%
  \Box!A \in \Delta \text{ se e somente se } !A \in \Delta' \text{
    para todo $\Delta'$ tal que } R^\Sigma\Delta\Delta'}{%
  \Diamond!A \in \Delta \text{ se e somente se } !A \in \Delta' \text{ para pelo
    menos um } \Delta' \text{ tal que } R^\Sigma\Delta\Delta'}
\]
\iftag{prvBox}{Considere a direção da esquerda para a direita: ela diz
  que, se $\Box !A \in \Delta$, então $!A \in \Delta'$ para qualquer
  $!A$ e qualquer $\Delta'$ com $R^\Sigma\Delta\Delta'$. Se estipularmos
  que $R^\Sigma\Delta\Delta'$ se e somente se $!A \in \Delta'$ para
  todo $\Box!A \in \Delta$, então isso vale. Podemos escrever a condição
  à direita de ``se e somente se'' de modo mais compacto como:
  $\Setabs{!A}{\Box!A \in \Delta} \subseteq \Delta'$.}{Considere a
  direção da direita para a esquerda: ela diz que, se $!A \in \Delta'$,
  então $\Diamond!A \in \Delta$, para qualquer $!A$ e $\Delta'$ com
  $R^\Sigma\Delta\Delta'$. Se estipularmos que
  $R^\Sigma\Delta\Delta'$ vale sempre que $\Diamond!A \in \Delta$ para
  todo $!A \in \Delta'$, então isso vale. Podemos escrever a condição
  à direita de ``se e somente se'' de modo mais compacto como:
  $\Setabs{\Diamond!A}{!A \in \Delta'} \subseteq \Delta$.}

Assim, a pergunta é: essa definição de $R^\Sigma$ garante de fato que
\iftag{prvBox}{$\Box !A \in \Delta$ se e somente se
  $\mSat{M^\Sigma}{\Box !A}[\Delta]$? \iftag{prvDiamond}{Ela também
    garante que }{}}{}\iftag{prvDiamond}{$\Diamond !A \in \Delta$ se e somente se
  $\mSat{M^\Sigma}{\Diamond !A}[\Delta]$?}{} Os próximos resultados
estabelecerão isso.
\end{explain}

\begin{defn}
  Se $\Gamma$ é um conjunto de !!{formula}s, seja
  \begin{align*}
    \Box\Gamma & = \Setabs{\Box !B}{!B \in \Gamma}\\
    \Diamond\Gamma & = \Setabs{\Diamond !B}{!B \in \Gamma}\\
    \intertext{e}
    \Box^{-1}\Gamma & = \Setabs{!B}{\Box !B \in \Gamma}\\
    \Diamond^{-1}\Gamma & = \Setabs{!B}{\Diamond !B \in \Gamma}\\
  \end{align*}
\end{defn}

Em outras palavras, $\Box\Gamma$ é $\Gamma$ com $\Box$ diante de cada
!!{formula} em~$\Gamma$; $\Box^{-1}\Gamma$ consiste em todas as
!!{formula}s com prefixo~$\Box$ que pertencem a~$\Gamma$, com os
símbolos $\Box$ iniciais removidos. Essa definição não é muito importante por si só,
mas simplificará consideravelmente a notação.

Observe que $\Box\Box^{-1}\Gamma \subseteq \Gamma$:
\[
\Box\Box^{-1}\Gamma = \Setabs{\Box!B}{\Box!B \in \Gamma}
\]
O.s.\@, trata-se apenas do conjunto de todas as !!{formula}s
de~$\Gamma$ que começam com~$\Box$.

\begin{lem}\ollabel{lem:box1}
  Se $\Gamma \Proves[\Sigma] !A$, então $\Box\Gamma \Proves[\Sigma] \Box
  !A$.
\end{lem}

\begin{proof}
  Se $\Gamma \Proves[\Sigma] !A$, então existem $!B_1$, \dots, $!B_k
  \in \Gamma$ tais que $\Sigma \Proves !B_1 \lif (!B_2 \lif \cdots
  (!B_n \lif !A)\cdots)$. Como $\Sigma$ é normal, pela regra \RK{},
  $\Sigma \Proves \Box!B_1 \lif (\Box!B_2 \lif \cdots (\Box!B_n \lif
  \Box!A)\cdots)$, onde, evidentemente, $\Box!B_1$, \dots,
  $\Box!B_k \in \Box\Gamma$. Logo, por definição,
  $\Box\Gamma \Proves[\Sigma] \Box
  !A$.
\end{proof}

\begin{lem}\ollabel{lem:box2}
  Se $\Box^{-1} \Gamma \Proves[\Sigma] !A$, então $\Gamma
  \Proves[\Sigma] \Box!A$.
\end{lem}

\begin{proof}
  Suponha que $\Box^{-1}\Gamma \Proves[\Sigma] !A$; então, por
  \olref{lem:box1}, $\Box\Box^{-1}\Gamma \Proves \Box!A$. Mas, como
  $\Box\Box^{-1}\Gamma \subseteq \Gamma$, também
  $\Gamma \Proves[\Sigma] \Box!A$ por monotonicidade.
\end{proof}

\iftag{prvBox}{%
% Box é primitivo
  
\begin{prop}\ollabel{prop:box}
  Se $\Gamma$ é $\Sigma$-consistente completo, então $\Box !A \in
  \Gamma$ se e somente se, para todo conjunto $\Sigma$-consistente
  completo $\Delta$ tal que $\Box^{-1} \Gamma \subseteq \Delta$, vale
  que $!A \in \Delta$.
\end{prop}

\begin{proof}
  Suponha que $\Gamma$ seja $\Sigma$-consistente completo. A direção
  ``somente se'' é fácil: suponha que $\Box !A \in \Gamma$ e que
  $\Box^{-1}\Gamma \subseteq \Delta$. Como $\Box!A \in \Gamma$,
  $!A \in \Box^{-1}\Gamma \subseteq \Delta$, logo $!A \in \Delta$.

  Para a direção ``se'', demonstramos a contrapositiva: suponha que
  $\Box!A \notin \Gamma$. Como $\Gamma$ é $\Sigma$-consistente
  completo, é dedutivamente fechado e, portanto, $\Gamma
  \Proves/[\Sigma] \Box !A$. Por \olref{lem:box2},
  $\Box^{-1}\Gamma \Proves/[\Sigma] !A$. Por
  \olref[prf][con]{prop:consistencyfacts}\olref[prf][con]{prop:consistencyfacts-b},
  $\Box^{-1}\Gamma \cup \{ \lnot!A \}$ é $\Sigma$-consistente. Pelo
  lema de Lindenbaum, existe um conjunto $\Sigma$-consistente completo
  $\Delta$ tal que $\Box^{-1}\Gamma \cup \{ \lnot!A \} \subseteq
  \Delta$. Pela consistência, $!A \notin \Delta$.
\end{proof}
}{%
% Box não é primitivo; apenas Diamond é

\begin{prop}\ollabel{prop:diamond}
  Se $\Gamma$ é $\Sigma$-consistente completo, então $\Diamond !A \in
  \Gamma$ se e somente se, para algum conjunto $\Sigma$-consistente
  completo $\Delta$ tal que $\Diamond\Delta \subseteq
  \Gamma$, vale que $!A \in \Delta$.
\end{prop}

\begin{proof}
  Suponha que $\Gamma$ seja $\Sigma$-consistente completo. A parte da
  direita para a esquerda é fácil: seja $\Delta$ $\Sigma$-consistente
  completo tal que $\Diamond\Delta \subseteq \Gamma$ e
  $!A \in \Delta$. Então $\Diamond!A \in \Diamond\Delta \subseteq
  \Gamma$, O.s.\@, $\Diamond!A \in \Gamma$.

  Para a direção da esquerda para a direita, suponha que
  $\Diamond !A \in \Gamma$. Devemos mostrar que existe um conjunto
  $\Sigma$-consistente completo $\Delta$ com $!A \in \Delta$ e
  $\Diamond\Delta \subseteq \Gamma$. Como $\Diamond !A \in \Gamma$ e
  $\Gamma$ é $\Sigma$-consistente, $\Box\lnot !A \notin \Gamma$. Como
  $\Gamma$ é $\Sigma$-consistente completo, é dedutivamente fechado,
  logo $\Gamma \Proves/[\Sigma] \Box\lnot !A$. Por \olref{lem:box2},
  $\Box^{-1}\Gamma \Proves/[\Sigma] \lnot !A$. Por
  \olref[prf][con]{prop:consistencyfacts}\olref[prf][con]{prop:consistencyfacts-b},
  $\Box^{-1}\Gamma \cup \{ !A \}$ é $\Sigma$-consistente. Pelo lema de
  Lindenbaum, existe um conjunto $\Sigma$-consistente completo
  $\Delta$ tal que $\Box^{-1}\Gamma \cup \{ !A \} \subseteq \Delta$.
  Claramente, $!A \in \Delta$. Para ver que $\Diamond\Delta
  \subseteq \Gamma$, suponha o contrário: para algum $!B \in \Delta$,
  $\Diamond !B \notin \Gamma$. Então, como $\Gamma$ é
  $\Sigma$-consistente completo, $\Box\lnot !B \in \Gamma$. Mas então
  também $\lnot !B \in \Box^{-1}\Gamma \subseteq \Delta$, o que
  tornaria $\Delta$ inconsistente.
\end{proof}
}

\begin{lem}\ollabel{lem:box-iff-diamond}
  Suponha que $\Gamma$ e $\Delta$ sejam $\Sigma$-consistentes
  completos. Então $\Box^{-1}\Gamma \subseteq \Delta$ se e somente se
  $\Diamond\Delta \subseteq \Gamma$.
\end{lem}

\begin{proof}
  Direção ``somente se'': suponha $\Box^{-1}\Gamma \subseteq \Delta$ e
  $\Diamond!A \in \Diamond\Delta$ (O.s.\@, $!A \in \Delta$). Para
  mostrar $\Diamond!A \in \Gamma$, basta mostrar
  $\Box\lnot!A \notin \Gamma$, pois então, pela maximalidade,
  $\lnot\Box\lnot !A \in \Gamma$. Agora, se
  $\Box\lnot!A \in \Gamma$, então, pela hipótese, $\lnot!A \in \Delta$,
  contrariando a consistência de $\Delta$ (pois $!A \in \Delta$).
  Logo, $\Box\lnot!A \notin \Gamma$, como exigido.

  Direção ``se'': suponha $\Diamond\Delta \subseteq \Gamma$.
  Argumentamos por contraposição: suponha $!A \notin \Delta$ para
  mostrar $\Box!A \notin \Gamma$. Se $!A \notin \Delta$, então, pela
  maximalidade, $\lnot!A \in \Delta$ e, portanto, pela hipótese,
  $\Diamond\lnot!A \in \Gamma$. Mas, em uma lógica modal normal,
  $\Diamond\lnot!A$ é equivalente a $\lnot\Box !A$ e, se esta última
  pertence a $\Gamma$, pela consistência $\Box!A \notin\Gamma$, como
  exigido.
\end{proof}

\iftag{prvBox}{%
  \iftag{prvDiamond}{%
% Box e Diamond são ambos primitivos; demonstre
% prop:diamond usando lem:box-iff-diamond
  
\begin{prop}\ollabel{prop:diamond}
  Se $\Gamma$ é $\Sigma$-consistente completo, então $\Diamond !A \in
  \Gamma$ se e somente se, para algum conjunto $\Sigma$-consistente
  completo $\Delta$ tal que $\Diamond\Delta \subseteq
  \Gamma$, vale que $!A \in \Delta$.
\end{prop}

\begin{proof}
Suponha que $\Gamma$ seja $\Sigma$-consistente completo.
$\Diamond!A \in \Gamma$ se e somente se
$\lnot\Box\lnot !A \in \Gamma$, por \Dual{} e pelo fechamento.
$\lnot\Box\lnot !A \in \Gamma$ se e somente se
$\Box\lnot !A \notin \Gamma$, por
\olref[ccs]{prop:ccs-properties}\olref[ccs]{prop:ccs-lnot}, pois
$\Gamma$ é $\Sigma$-consistente completo. Por \olref{prop:box},
$\Box\lnot !A \notin \Gamma$ se e somente se, para algum conjunto
$\Sigma$-consistente completo $\Delta$ com
$\Box^{-1}\Gamma \subseteq \Delta$, $\lnot!A \notin \Delta$.
Considere agora qualquer $\Delta$ desse tipo. Por
\olref{lem:box-iff-diamond}, $\Box^{-1}\Gamma \subseteq \Delta$ se e
somente se $\Diamond\Delta \subseteq \Gamma$. Além disso,
$\lnot !A \notin \Delta$ se e somente se $!A \in \Delta$, por
\olref[ccs]{prop:ccs-properties}\olref[ccs]{prop:ccs-lnot}. Portanto,
$\Diamond!A \in \Gamma$ se e somente se, para algum conjunto
$\Sigma$-consistente completo $\Delta$ com
$\Diamond\Delta \subseteq \Gamma$, $!A \in \Delta$.
\end{proof}

}{}}{}
  
\begin{prob}
  Mostre que, se $\Gamma$ é $\Sigma$-consistente completo, então
  \iftag{prvBox}{$\Diamond !A \in \Gamma$ se e somente se existe um
    conjunto $\Sigma$-consistente completo $\Delta$ tal que $\Box^{-1}\Gamma
    \subseteq \Delta$ e $!A \in \Delta$.}{$\Box !A \in
    \Gamma$ se e somente se, para todo conjunto $\Sigma$-consistente
    completo $\Delta$ tal que $\Box^{-1} \Gamma \subseteq \Delta$, vale
    que $!A \in \Delta$.} \emph{Faça isso sem usar
    \olref[nml][com][mod]{lem:box-iff-diamond}}.
\end{prob}

\end{document}
